% Generated by scripts/make_table_main.py. Never hand-edited.
% Every number traces to a cell recorded in results/table_main_cells.csv.
% Requires \usepackage{booktabs}.
% - Gate G3: PASS. It is a completeness gate and
%   it is not the claim; Table 3 is what the claim measured.
% - Pre-declared primary quantity, recovery\_ratio at
%   deterioration\_mortality\_365d: NOT MET.
% - All four pre-declared quantitative targets are in Table 3 and none was met
%   at the endpoint it named.

% PANEL T1
\begin{table}[t]
\centering
\small
\caption{The nested input ladder, record level. AUROC with a patient-clustered percentile bootstrap interval, 10,000 resamples, resampling unit subject\_id. One row per ECG throughout: rows 1 to 5 have no crops, so the published crop-level protocol appears nowhere in this panel. Scores are averaged across the five seeds and one AUROC is taken on the pooled score, because a paired bootstrap needs one score vector per arm; per-seed spreads are in Additional file 1, Table S7.}
\label{tab:ladder}
\begin{tabular}{lllll}
\toprule
Row & Arm & ICU admission, 24h & Mortality, 28d & Mortality, 365d \\
\midrule
1 & \texttt{R1\_prevalence} & 0.5000 [0.5000, 0.5000] & 0.5000 [0.5000, 0.5000] & 0.5000 [0.5000, 0.5000] \\
2 & \texttt{R2\_demo} & 0.6017 [0.5752, 0.6271] & 0.6963 [0.6608, 0.7307] & 0.7458 [0.7205, 0.7696] \\
3 & \texttt{R3\_acqctx\_pre} & 0.7972 [0.7806, 0.8138] & 0.7587 [0.7264, 0.7903] & 0.7241 [0.7029, 0.7441] \\
3b & \texttt{R3b\_acqctx\_window} & 0.8010 [0.7845, 0.8174] & 0.7637 [0.7303, 0.7958] & 0.7272 [0.7057, 0.7474] \\
4 & \texttt{R4\_demo\_acq} & 0.8094 [0.7928, 0.8255] & 0.8102 [0.7818, 0.8371] & 0.8033 [0.7832, 0.8221] \\
5 & \texttt{R5\_tabular} & 0.9164 [0.9046, 0.9278] & 0.9102 [0.8884, 0.9294] & 0.8533 [0.8371, 0.8689] \\
6 & \texttt{R6\_waveform} & 0.8173 [0.7995, 0.8342] & 0.8473 [0.8222, 0.8714] & 0.8294 [0.8109, 0.8474] \\
7 & \texttt{R7\_waveform\_demo\_acq} & 0.8571 [0.8424, 0.8713] & 0.8723 [0.8485, 0.8947] & 0.8533 [0.8365, 0.8693] \\
\bottomrule
\end{tabular}
\par\footnotesize Row 8 is not a row here. It is a within-stratum concordance over restricted pairs, a different estimand from a record-level AUROC, and it has its own panel in Additional file 1, Table S1.
\par\footnotesize Sample: ICU admission, 24h: 6,407 ECGs, 3,626 patients, prevalence 0.1260; Mortality, 28d: 6,428 ECGs, 3,636 patients, prevalence 0.0422; Mortality, 365d: 6,428 ECGs, 3,636 patients, prevalence 0.1425.
\end{table}

% PANEL T2
\begin{table}[t]
\centering
\small
\caption{Every pre-declared quantitative target, and what it measured. Each target was fixed in a committed comparison file before the run that tests it existed.}
\label{tab:targets}
\begin{tabular}{lllll}
\toprule
Quantity & Endpoint & Pre-declared & Measured [95\% CI] & Outcome \\
\midrule
pre-declared rule 1: R3\_acqctx\_pre beats R2\_demo & Mortality, 365d & positive, beyond the 0.0043 seed and noise spread & -0.0216 [-0.0523, +0.0098] & FAILED \\
Recovery ratio, R3\_acqctx\_pre against R6\_waveform & Mortality, 365d & at least 0.70 & 0.6803 [0.6149, 0.7476] & FAILED \\
Waveform gain retained with acquisition context held fixed & Mortality, 365d & at most 0.50 & 0.9313 [0.6631, 1.1583] & FAILED \\
Horizon-ladder trend, four never-scored horizons & mortality 1, 7, 90, 180 d & negative, interval excluding zero & -0.8000 [-1.0000, +0.2000] & NOT SUPPORTED \\
\bottomrule
\end{tabular}
\par\footnotesize The recovery ratio clears 0.70 at the two labels that were not the pre-declared endpoint (Table 4); both are reported as exploratory.
\par\footnotesize The retained-gain target is missed at all three labels and under both readings of its ambiguous formula (Table 6), not only at the endpoint shown here. Under the second reading the intervals at 28 days and at one year contain the target, so at those two the evidence is inconclusive rather than contrary; the outcome column reports the declared rule, which is read on the point estimate.
\end{table}

% PANEL T3
\begin{table}[t]
\centering
\small
\caption{Recovery ratio against the waveform arm. (AUROC\_arm - 0.5) / (AUROC\_R6 - 0.5), with the interval taken on the ratio itself: patients are resampled once and both AUROCs are recomputed on that same draw. The denominator is the per-label record-level R6\_waveform, not the 15-target macro; the two readings differ enough to move the primary across its own threshold.}
\label{tab:recovery}
\begin{tabular}{lllllll}
\toprule
Arm & Label & Ratio [95\% CI] & Numerator & Denominator & At least 0.70 & Status \\
\midrule
\texttt{R3\_acqctx\_pre} & ICU admission, 24h & 0.9366 [0.8758, 0.9995] & 0.7972 & 0.8173 & yes & exploratory \\
\texttt{R3\_acqctx\_pre} & Mortality, 28d & 0.7448 [0.6515, 0.8382] & 0.7587 & 0.8473 & yes & exploratory \\
\texttt{R3\_acqctx\_pre} & Mortality, 365d & 0.6803 [0.6149, 0.7476] & 0.7241 & 0.8294 & no & pre-declared primary \\
\texttt{R4\_demo\_acq} & ICU admission, 24h & 0.9749 [0.9155, 1.0376] & 0.8094 & 0.8173 & yes & exploratory \\
\texttt{R4\_demo\_acq} & Mortality, 28d & 0.8931 [0.8095, 0.9769] & 0.8102 & 0.8473 & yes & exploratory \\
\texttt{R4\_demo\_acq} & Mortality, 365d & 0.9206 [0.8568, 0.9849] & 0.8033 & 0.8294 & yes & exploratory \\
\texttt{R5\_tabular} & ICU admission, 24h & 1.3122 [1.2497, 1.3833] & 0.9164 & 0.8173 & yes & exploratory \\
\texttt{R5\_tabular} & Mortality, 28d & 1.1811 [1.1132, 1.2564] & 0.9102 & 0.8473 & yes & exploratory \\
\texttt{R5\_tabular} & Mortality, 365d & 1.0726 [1.0204, 1.1294] & 0.8533 & 0.8294 & yes & exploratory \\
\texttt{R7\_waveform\_demo\_acq} & ICU admission, 24h & 1.1253 [1.0913, 1.1624] & 0.8571 & 0.8173 & yes & exploratory \\
\texttt{R7\_waveform\_demo\_acq} & Mortality, 28d & 1.0720 [1.0357, 1.1119] & 0.8723 & 0.8473 & yes & exploratory \\
\texttt{R7\_waveform\_demo\_acq} & Mortality, 365d & 1.0725 [1.0457, 1.1010] & 0.8533 & 0.8294 & yes & exploratory \\
\bottomrule
\end{tabular}
\par\footnotesize R5\_tabular, the benchmark's own vitals-and-labs arm, exceeds 1.0 at every horizon: it beats the waveform arm it is bundled with. The benchmark's own appendix reports the same ordering, so this replicates a published direction at a per-label unit rather than reporting a new one.
\end{table}

% PANEL T3b
\begin{table}[t]
\centering
\small
\caption{The primary recovery ratio under every declared denominator. The registered primary quantity divides the acquisition-context arm's above-chance discrimination by the waveform arm's, at the registered endpoint. Which waveform figure belongs in the denominator was declared with two alternatives beside the one used, and a third became available when the audited benchmark's version of record revised its published macro. All four are here. Only the declared denominator is resampled with the numerator. A published macro is a constant from another paper, and our own 15-target macro is a different estimand from the per-label AUROC, so for those three the interval carries the numerator's uncertainty alone and the table says which.}
\label{tab:recovery_denominators}
\begin{tabular}{lllll}
\toprule
Denominator & Value & Ratio [95\% CI] & Resampled & At least 0.70 \\
\midrule
per-label `R6\_waveform` (declared) & 0.8294 & 0.6803 [0.6149, 0.7476] & both & no \\
published macro, version of record & 0.8565 & 0.6286 [0.5692, 0.6847] & numerator only & no \\
published macro, preprint & 0.8279 & 0.6835 [0.6189, 0.7445] & numerator only & no \\
our own 15-target record-level macro & 0.8545 & 0.6322 [0.5724, 0.6886] & numerator only & no \\
\bottomrule
\end{tabular}
\par\footnotesize No denominator brings the ratio to its registered target. The declared one is the most favourable of the four: under it the interval spans 0.70, while under the version-of-record macro and under our own 15-target macro the whole interval falls below the target. The preprint macro gives a slightly larger ratio than the declared denominator rather than a smaller one, which is the opposite of what this paper asserted before the cells existed.
\end{table}

% PANEL T0
\begin{table}[t]
\centering
\small
\caption{The reproduced cohort against the published counts. Every count the release states, beside ours, with the difference. Patients and stays match to the unit and the record count falls 38 short, which is the shape a partially withheld release leaves and not the shape a differently built cohort would. The sequence number within a stay is ZERO-based here: a record numbered 0 is the first ECG of its stay. The benchmark's own instructions restrict validation and test scoring to that first record and this audit scores every test record, so the two use different denominators and this table gives both.}
\label{tab:cohort_reconciliation}
\begin{tabular}{lllll}
\toprule
Section & Quantity & Local & Published & Gap \\
\midrule
Cohort & ECG records & 129,057 & 129,095 & -38 \\
Cohort & Patients & 71,098 & 71,098 & +0 \\
Cohort & Emergency department stays & 121,195 & 121,195 & +0 \\
Label positives & deterioration\_icu\_24h & 15,677 & 15,688 & -11 \\
Label exclusions & deterioration\_icu\_24h & 539 & not published & - \\
Label positives & deterioration\_icu\_stay & 18,860 & 18,872 & -12 \\
Label exclusions & deterioration\_icu\_stay & 539 & not published & - \\
Label positives & deterioration\_mortality\_28d & 5,150 & 5,153 & -3 \\
Label exclusions & deterioration\_mortality\_28d & 357 & not published & - \\
Label positives & deterioration\_mortality\_365d & 17,756 & 17,761 & -5 \\
Label exclusions & deterioration\_mortality\_365d & 357 & not published & - \\
Split & Training folds 0 to 17 & 116,433 & not published & - \\
Split & Validation fold 18 & 6,185 & not published & - \\
Split & Test fold 19 & 6,439 & not published & - \\
Test fold & ECG records & 6,439 & not published & - \\
Test fold & Patients & 3,644 & not published & - \\
Test fold & Stays & 6,077 & not published & - \\
First ECG per stay & Records numbered 0 within their stay & 6,080 & not published & - \\
First ECG per stay & Distinct stays among them & 6,076 & not published & - \\
First ECG per stay & Stays carrying two records numbered 0 & 4 & not published & - \\
First ECG per stay & Records after the first of their stay & 359 & not published & - \\
Sequence within stay & Records numbered 0 & 6,080 & not published & - \\
Sequence within stay & Records numbered 1 & 318 & not published & - \\
Sequence within stay & Records numbered 2 & 33 & not published & - \\
Sequence within stay & Records numbered 3 & 6 & not published & - \\
Sequence within stay & Records numbered 4 & 2 & not published & - \\
\bottomrule
\end{tabular}
\par\footnotesize Label exclusions are records whose outcome is not observable for that label, marked in the release with a sentinel value. The tabular arms drop them, which is why the analysed sample differs slightly by label.
\par\footnotesize Four stays carry two records both numbered 0, because their acquisition times tie and the ordering does not break ties. The first ECG of those stays is therefore ambiguous, and the first-record restriction covers four fewer distinct stays than it has records.
\end{table}

% PANEL T8b
\begin{table}[t]
\centering
\small
\caption{The waveform reproduction under both scoring protocols. EXPLORATORY, and not registered. The audited benchmark scores validation and test on the first record of each visit; this audit scores every test record, so the two use different denominators. Restricting to the first record is a selection over predictions that already exist, so nothing is retrained and the arm is the one every other table reports. The per-seed range is beside each macro, because a difference smaller than the spread of the five seeds is not a difference.}
\label{tab:first_ecg}
\begin{tabular}{lllll}
\toprule
Unit & Scoring protocol & Rows scored & Macro AUROC & Per-seed range \\
\midrule
crop & all test records & 25,756 & 0.8474 & 0.0087 \\
crop & first ECG of each stay & 24,320 & 0.8465 & 0.0097 \\
record & all test records & 6,439 & 0.8545 & 0.0102 \\
record & first ECG of each stay & 6,080 & 0.8536 & 0.0111 \\
\bottomrule
\end{tabular}
\par\footnotesize The restriction moves the macro by about a thousandth at both units, an order of magnitude below the seed-to-seed range, and the restricted macro falls inside the published interval exactly as the unrestricted one does. The difference between the two protocols is therefore not what separates this reproduction from the published figure.
\par\footnotesize The restriction keeps 6,080 records over 6,076 stays: four stays carry two records whose acquisition times tie, and both are kept rather than resolved by a tie-break nobody declared.
\end{table}

% PANEL T1b
\begin{table}[t]
\centering
\small
\caption{The acquisition-context block against triage acuity alone. EXPLORATORY, and not registered. Both reviews asked how much the acquisition-context block adds beyond the clinical acuity a triage nurse already records, and no arm in the paper isolated that feature. This one does: a single ordinal column, fitted the same way as every other tabular arm. The last row of each label is the paired difference between the nine-feature arm and the one-feature arm, on the same records and the same resampled patients, which is the quantity the question is actually about.}
\label{tab:triage}
\begin{tabular}{lll}
\toprule
Label & Arm & AUROC or difference [95\% CI] \\
\midrule
ICU admission, 24h & `R9\_triage`, triage acuity alone & 0.7583 [0.7410, 0.7755] \\
ICU admission, 24h & `R2\_demo`, age and sex & 0.6017 [0.5752, 0.6271] \\
ICU admission, 24h & `R3\_acqctx\_pre`, all nine features & 0.7972 [0.7806, 0.8138] \\
ICU admission, 24h & what the other eight features add & +0.0389 [+0.0310, +0.0467] \\
Mortality, 28d & `R9\_triage`, triage acuity alone & 0.7120 [0.6757, 0.7480] \\
Mortality, 28d & `R2\_demo`, age and sex & 0.6963 [0.6608, 0.7307] \\
Mortality, 28d & `R3\_acqctx\_pre`, all nine features & 0.7587 [0.7264, 0.7903] \\
Mortality, 28d & what the other eight features add & +0.0467 [+0.0266, +0.0675] \\
Mortality, 365d & `R9\_triage`, triage acuity alone & 0.6298 [0.6095, 0.6502] \\
Mortality, 365d & `R2\_demo`, age and sex & 0.7458 [0.7205, 0.7696] \\
Mortality, 365d & `R3\_acqctx\_pre`, all nine features & 0.7241 [0.7029, 0.7441] \\
Mortality, 365d & what the other eight features add & +0.0943 [+0.0736, +0.1146] \\
\bottomrule
\end{tabular}
\par\footnotesize Acuity alone carries most of the block's discrimination at 24-hour ICU admission and the other eight features still add to it, at every label, with every interval excluding zero. The block is neither reducible to acuity nor independent of it.
\par\footnotesize At one year the ordering between acuity and demographics reverses: acuity alone falls below age and sex, having been well above them at 24 hours. That is consistent with the intuition the horizon-ladder family was built on, and it is not evidence for that family's registered trend, which the paper reports as not supported. An exploratory arm cannot rescue a registered test.
\par\footnotesize This arm fits one ordinal feature, so every seed produces the same model and its seed-to-seed spread is exactly zero. The noise floor this paper judges differences against therefore does not exist for it, and the interval carries the uncertainty instead.
\end{table}

% PANEL T4f
\begin{table}[t]
\centering
\small
\caption{The declared matching against a sensitivity without the weekday. EXPLORATORY, and it moves no reported result. The declared specification strata on four coarsened variables and row 8 is computed under it. One of those four is the recorded day of the week, which this dataset's de-identification leaves carrying no information, so the split costs pairs and controls nothing. This is that cost, measured by re-cutting the strata on the three informative variables and counting what the within-stratum statistic can then use.}
\label{tab:matching_sensitivity}
\begin{tabular}{lllll}
\toprule
Label & Specification & Strata & Pairs retained & Effective patients \\
\midrule
ICU admission, 24h & Declared & 143 & 0.0157 & 3,548 \\
ICU admission, 24h & Without weekday & 76 & 0.0306 & 3,615 \\
Mortality, 28d & Declared & 143 & 0.0155 & 3,287 \\
Mortality, 28d & Without weekday & 76 & 0.0299 & 3,507 \\
Mortality, 365d & Declared & 143 & 0.0199 & 3,582 \\
Mortality, 365d & Without weekday & 76 & 0.0388 & 3,615 \\
\bottomrule
\end{tabular}
\par\footnotesize Dropping the weekday roughly halves the number of strata and roughly doubles the share of pairs the statistic can compare, while leaving the same patients defined. The declared analysis is not re-cut: a specification fixed before the numbers existed is not revised once they do, and what this supports is a statement about the precision of row 8 rather than about its estimate.
\end{table}

% PANEL T4a
\begin{table}[t]
\centering
\small
\caption{Row 8: discrimination with acquisition context held fixed. A different estimand from Table 3, and never in the same column as it. Only positive-negative pairs drawn from the same coarsened acquisition-context stratum contribute, so no compared pair differs in acquisition context. The plain column is computed on the same rows, so the only difference between the two is the pair restriction rather than which records were kept.}
\label{tab:row8}
\begin{tabular}{llllll}
\toprule
Arm & Label & Within stratum [95\% CI] & Plain, same rows & Effective patients & Pairs retained \\
\midrule
\texttt{R2\_demo} & ICU admission, 24h & 0.5796 [0.5460, 0.6144] & 0.6026 & 3,548 & 0.0157 \\
\texttt{R3\_acqctx\_pre} & ICU admission, 24h & 0.6042 [0.5735, 0.6349] & 0.7888 & 3,548 & 0.0157 \\
\texttt{R8\_waveform\_matched} & ICU admission, 24h & 0.7464 [0.7149, 0.7758] & 0.8152 & 3,548 & 0.0157 \\
\texttt{R2\_demo} & Mortality, 28d & 0.6464 [0.5922, 0.6985] & 0.6938 & 3,287 & 0.0155 \\
\texttt{R3\_acqctx\_pre} & Mortality, 28d & 0.6158 [0.5595, 0.6702] & 0.7554 & 3,287 & 0.0155 \\
\texttt{R8\_waveform\_matched} & Mortality, 28d & 0.7776 [0.7198, 0.8324] & 0.8456 & 3,287 & 0.0155 \\
\texttt{R2\_demo} & Mortality, 365d & 0.7215 [0.6903, 0.7537] & 0.7476 & 3,582 & 0.0199 \\
\texttt{R3\_acqctx\_pre} & Mortality, 365d & 0.6792 [0.6490, 0.7083] & 0.7236 & 3,582 & 0.0199 \\
\texttt{R8\_waveform\_matched} & Mortality, 365d & 0.7984 [0.7708, 0.8250] & 0.8301 & 3,582 & 0.0199 \\
\bottomrule
\end{tabular}
\par\footnotesize R3\_acqctx\_pre collapses within a stratum because four of its nine features are the matching variables, R2\_demo barely moves, and the waveform arm keeps most of its advantage. A statistic that failed to restrict its pairs could not produce that pattern.
\par\footnotesize The effective patient count, not the ECG count, bounds every claim in this panel. Almost every patient survives the restriction and most pairs do not, which is why these intervals are two to three times wider than those of Table 3 on the same records.
\end{table}

% PANEL T4b
\begin{table}[t]
\centering
\small
\caption{Waveform gain retained: the declared secondary quantity. (R6 - R2) within stratum over (R6 - R2) on the same rows without the pair restriction. The declared formula does not say which R2 the numerator uses; both readings are reported and neither was chosen and they agree on the verdict.}
\label{tab:retained}
\begin{tabular}{llllll}
\toprule
Label & Retained [95\% CI] & Meets target & Retained, unstratified R2 [95\% CI] & Meets target & Declared target \\
\midrule
ICU admission, 24h & 0.7846 [0.6308, 0.9259] & no & 0.6765 [0.5491, 0.7830] & no & at most 0.50 \\
Mortality, 28d & 0.8643 [0.5584, 1.1720] & no & 0.5518 [0.1822, 0.8262] & no & at most 0.50 \\
Mortality, 365d & 0.9313 [0.6631, 1.1583] & no & 0.6148 [0.3014, 0.8244] & no & at most 0.50 \\
\bottomrule
\end{tabular}
\par\footnotesize The prediction was that holding acquisition context fixed would cost the waveform arm at least half its advantage over demographics. Under the within-stratum reading of the numerator it costs between 7 and 22 percent of it; under the unstratified reading, between 32 and 45 percent. The point estimate misses the declared target under both. Under the second reading the intervals at 28 days and at one year include 0.50, so at those two labels the evidence does not settle which side of the target the truth lies on; what is established is that the target is not met under the declared rule, which is read on the point estimate.
\end{table}

% PANEL T4c
\begin{table}[t]
\centering
\small
\caption{Balance, and the acquisition-propensity diagnostic. Standardised mean differences between outcome-positive and outcome-negative records, under the pair weighting the within-stratum statistic implies. On the coarsened matching variables balance is exact after weighting by construction, so a non-zero value there would be a bug in the stratifier rather than a property of the data; a shuffled stratum assignment breaks those zeros, so they are not the output of an insensitive statistic. The residual on the underlying continuous values is what coarsening leaves inside a bin and is reported rather than assumed away.}
\label{tab:balance}
\begin{tabular}{llll}
\toprule
Scale & Worst |SMD| before & Worst |SMD| after & Comparisons \\
\midrule
matching variables, as coarsened & 0.7889 & 0.0000 & 45 \\
matching variables, underlying scale & 1.0013 & 0.0993 & 12 \\
pre-acquisition features the strata do not hold fixed & 0.3711 & 0.3881 & 15 \\
\bottomrule
\end{tabular}
\par\footnotesize As a diagnostic only, pre-acquisition context predicts whether an early ECG happens at all with AUROC 0.6764 on held-out stays from the parent cohort, where the early-ECG rate is 0.2851. So the decision to record is itself substantially predictable. That model is never used to select or weight a row 8 evaluation row: the parent cohort carries no deterioration label and its ECG-absent stays carry no waveform.
\end{table}

% PANEL T4d
\begin{table}[t]
\centering
\small
\caption{Every declared acquisition-context feature, and what it can carry. Level counts over all 129,057 records of the reproduced cohort. A feature with one level carries no information, so it cannot be a leak whatever its availability stamp says, and it cannot be part of a confirmatory arm's discrimination either. Both directions matter here: the ACQ\_WINDOW block is where the leakage objection lives, and the ACQ\_PRE block is what the confirmatory claim rests on.}
\label{tab:blocks}
\begin{tabular}{lllllll}
\toprule
Feature & Block & Status & Carries information & Levels & Modal share & Undefined share \\
\midrule
\texttt{`arrival\_to\_acquisition\_minutes`} & ACQ\_PRE & varies & yes & 90 & 0.0702 & 0.0000 \\
\texttt{`acquisition\_hour\_of\_day`} & ACQ\_PRE & varies & yes & 24 & 0.0728 & 0.0000 \\
\texttt{`acquisition\_day\_of\_week`} & ACQ\_PRE & varies & yes & 7 & 0.1450 & 0.0000 \\
\texttt{`acquisition\_month`} & ACQ\_PRE & varies & yes & 12 & 0.0861 & 0.0000 \\
\texttt{`triage\_acuity`} & ACQ\_PRE & varies & yes & 5 & 0.5198 & 0.0233 \\
\texttt{`ecg\_no\_within\_stay`} & ACQ\_PRE & varies & yes & 7 & 0.9397 & 0.0000 \\
\texttt{`prior\_ecg\_exists`} & ACQ\_PRE & varies & yes & 2 & 0.6789 & 0.0000 \\
\texttt{`days\_since\_prior\_ecg`} & ACQ\_PRE & varies & yes & 70207 & 0.0064 & 0.3211 \\
\texttt{`prior\_ed\_visit\_count`} & ACQ\_PRE & varies & yes & 154 & 0.4468 & 0.0000 \\
\texttt{`n\_ecgs\_first\_90min`} & ACQ\_WINDOW & varies & yes & 7 & 0.8867 & 0.0000 \\
\texttt{`inter\_ecg\_interval\_minutes`} & ACQ\_WINDOW & varies & yes & 189 & 0.0605 & 0.8867 \\
\texttt{`care\_unit\_at\_acquisition`} & ACQ\_WINDOW & constant & no & 1 & 1.0000 & 0.0000 \\
\texttt{`cardiologist\_overread\_exists`} & ACQ\_WINDOW & dropped & no & - & - & - \\
\bottomrule
\end{tabular}
\par\footnotesize care\_unit\_at\_acquisition is constant because the benchmark admits only patients with an ECG inside the first 90 minutes of the ED stay, so the patient is in the Emergency Department at acquisition by construction. The benchmark's own inclusion rule forecloses the leak the ACQ\_WINDOW block was designed to isolate.
\par\footnotesize cardiologist\_overread\_exists was declared in the block, probed at data access, and dropped: MIMIC-IV-ECG ships machine-generated statements and carries no separate human over-read marker. It entered no run.
\par\footnotesize Every ACQ\_PRE feature varies. The confirmatory arm is fitted on nine features and all nine carry information.
\end{table}

% PANEL T4e
\begin{table}[t]
\centering
\small
\caption{Controls on the within-stratum statistic. Computed on synthetic data with a fixed seed, so these are properties of the estimator rather than results about the cohort. The third row is the one that matters: a score that is a pure function of the stratum discriminates well unstratified and must collapse to exactly chance once pairs are restricted to within a stratum. A statistic that did not really restrict its pairs would keep most of it.}
\label{tab:control}
\begin{tabular}{llll}
\toprule
Control & Unstratified & Within stratum & What it shows \\
\midrule
single stratum & 0.7384 & 0.7384 & one stratum reproduces plain AUROC exactly \\
constant within stratum & 0.7406 & 0.5000 & a score constant within every stratum collapses to chance \\
duplicated rows & 0.7384 & 0.7384 & duplicating every row leaves the statistic unchanged \\
\bottomrule
\end{tabular}
\par\footnotesize stratified.control\_values() is the single definition of these numbers. The module's own selftest asserts against it and this table is written from it, so the assertion and the printed figure cannot diverge.
\end{table}

% PANEL T5
\begin{table}[t]
\centering
\small
\caption{The audit family, complete at all 15 declared tests. acq\_context\_audit declared 5 contrasts crossed with 3 labels and the family has not grown. This is the correction the manuscript quotes: completing the family can only lower an existing adjusted p-value, never raise one, so the earlier correction is superseded here rather than contradicted.}
\label{tab:family}
\begin{tabular}{llllll}
\toprule
Contrast & Label & Raw p & Holm p, final & Phase & Family size m \\
\midrule
\texttt{R3\_acqctx\_pre - R2\_demo} & ICU admission, 24h & <0.0002 & <0.0002 & P3 & 15 \\
\texttt{R3\_acqctx\_pre - R6\_waveform} & Mortality, 28d & <0.0002 & <0.0002 & P3 & 15 \\
\texttt{R3\_acqctx\_pre - R6\_waveform} & Mortality, 365d & <0.0002 & <0.0002 & P3 & 15 \\
\texttt{R7\_waveform\_demo\_acq - R6\_waveform} & ICU admission, 24h & <0.0002 & <0.0002 & P3 & 15 \\
\texttt{R7\_waveform\_demo\_acq - R6\_waveform} & Mortality, 28d & <0.0002 & <0.0002 & P3 & 15 \\
\texttt{R7\_waveform\_demo\_acq - R6\_waveform} & Mortality, 365d & <0.0002 & <0.0002 & P3 & 15 \\
\texttt{R8\_waveform\_matched - R6\_waveform} & ICU admission, 24h & <0.0002 & <0.0002 & P4 & 15 \\
\texttt{R8\_waveform\_matched - R6\_waveform} & Mortality, 28d & <0.0002 & <0.0002 & P4 & 15 \\
\texttt{R8\_waveform\_matched - R6\_waveform} & Mortality, 365d & <0.0002 & <0.0002 & P4 & 15 \\
\texttt{R3\_acqctx\_pre - R2\_demo} & Mortality, 28d & 0.0108 & 0.0648 & P3 & 15 \\
\texttt{R4\_demo\_acq - R6\_waveform} & Mortality, 28d & 0.0128 & 0.0648 & P3 & 15 \\
\texttt{R4\_demo\_acq - R6\_waveform} & Mortality, 365d & 0.0152 & 0.0648 & P3 & 15 \\
\texttt{R3\_acqctx\_pre - R6\_waveform} & ICU admission, 24h & 0.0478 & 0.1434 & P3 & 15 \\
\texttt{R3\_acqctx\_pre - R2\_demo} & Mortality, 365d & 0.177 & 0.354 & P3 & 15 \\
\texttt{R4\_demo\_acq - R6\_waveform} & ICU admission, 24h & 0.4368 & 0.4368 & P3 & 15 \\
\bottomrule
\end{tabular}
\par\footnotesize This family is never pooled with the horizon ladder (Table 5), which carries its own Holm correction at its own declared size of 6.
\end{table}

% PANEL T5b
\begin{table}[t]
\centering
\small
\caption{Every declared contrast, as a paired difference. Paired patient-clustered bootstrap on the difference itself: both arms are scored on the same resampled patients in each replicate, and two marginal intervals are never compared. No p-value appears without its Each row's Holm value is from the completed family of 15 , not from the correction that was true when its phase closed. The noise floor is the seed-to-seed spread of that contrast's own arms at that label, and a difference not clearly larger than it is reported as no difference.}
\label{tab:contrasts}
\begin{tabular}{lllllll}
\toprule
Contrast & Label & Difference [95\% CI] & Holm p & Noise floor & Clears floor & Phase \\
\midrule
\texttt{R3\_acqctx\_pre - R2\_demo} & ICU admission, 24h & +0.1955 [+0.1649, +0.2267] & <0.0002 & 0.0049 & yes & P3 \\
\texttt{R3\_acqctx\_pre - R6\_waveform} & Mortality, 28d & -0.0886 [-0.1246, -0.0544] & <0.0002 & 0.0207 & yes & P3 \\
\texttt{R3\_acqctx\_pre - R6\_waveform} & Mortality, 365d & -0.1053 [-0.1301, -0.0809] & <0.0002 & 0.0130 & yes & P3 \\
\texttt{R7\_waveform\_demo\_acq - R6\_waveform} & ICU admission, 24h & +0.0398 [+0.0298, +0.0498] & <0.0002 & 0.0091 & yes & P3 \\
\texttt{R7\_waveform\_demo\_acq - R6\_waveform} & Mortality, 28d & +0.0250 [+0.0128, +0.0374] & <0.0002 & 0.0190 & yes & P3 \\
\texttt{R7\_waveform\_demo\_acq - R6\_waveform} & Mortality, 365d & +0.0239 [+0.0155, +0.0324] & <0.0002 & 0.0130 & yes & P3 \\
\texttt{R8\_waveform\_matched - R6\_waveform} & ICU admission, 24h & -0.0688 [-0.0912, -0.0478] & <0.0002 & - & - & P4 \\
\texttt{R8\_waveform\_matched - R6\_waveform} & Mortality, 28d & -0.0680 [-0.1102, -0.0284] & <0.0002 & - & - & P4 \\
\texttt{R8\_waveform\_matched - R6\_waveform} & Mortality, 365d & -0.0318 [-0.0488, -0.0157] & <0.0002 & - & - & P4 \\
\texttt{R3\_acqctx\_pre - R2\_demo} & Mortality, 28d & +0.0624 [+0.0165, +0.1082] & 0.0648 & 0.0207 & yes & P3 \\
\texttt{R4\_demo\_acq - R6\_waveform} & Mortality, 28d & -0.0371 [-0.0684, -0.0077] & 0.0648 & 0.0190 & yes & P3 \\
\texttt{R4\_demo\_acq - R6\_waveform} & Mortality, 365d & -0.0261 [-0.0481, -0.0049] & 0.0648 & 0.0130 & yes & P3 \\
\texttt{R3\_acqctx\_pre - R6\_waveform} & ICU admission, 24h & -0.0201 [-0.0406, -0.0002] & 0.1434 & 0.0091 & yes & P3 \\
\texttt{R3\_acqctx\_pre - R2\_demo} & Mortality, 365d & -0.0216 [-0.0523, +0.0098] & 0.354 & 0.0043 & yes & P3 \\
\texttt{R4\_demo\_acq - R6\_waveform} & ICU admission, 24h & -0.0080 [-0.0276, +0.0115] & 0.4368 & 0.0091 & no & P3 \\
\bottomrule
\end{tabular}
\par\footnotesize R4\_demo\_acq - R6\_waveform at 24-hour ICU admission does not clear its own noise floor. Demographics plus acquisition context is not distinguishable from the waveform arm there.
\par\footnotesize The three row 8 rows compare a within-stratum concordance against a plain AUROC on the same records, so they measure what the pair restriction costs rather than a difference between two models.
\end{table}

% PANEL T6a
\begin{table}[t]
\centering
\small
\caption{The horizon ladder: a separate family, per-horizon differences. R3\_acqctx\_pre minus R2\_demo at each mortality horizon. The primary test covers only the four horizons no arm had been scored on. The two horizons that generated the hypothesis are shown, carry no ladder-adjusted p, and never count as confirmation.}
\label{tab:ladder_horizons}
\begin{tabular}{llllll}
\toprule
Label & Horizon (d) & Difference [95\% CI] & Raw p & Holm p & Role \\
\midrule
Mortality, 1d & 1 & +0.0946 [-0.0360, +0.2168] & 0.1478 & 0.708 & primary \\
Mortality, 7d & 7 & +0.1224 [+0.0436, +0.2004] & 0.0006 & 0.0036 & primary \\
Mortality, 28d & 28 & +0.0624 [+0.0165, +0.1082] & 0.0108 & not corrected in this family & generating \\
Mortality, 90d & 90 & +0.0161 [-0.0225, +0.0539] & 0.4092 & 1 & primary \\
Mortality, 180d & 180 & +0.0030 [-0.0306, +0.0367] & 0.8656 & 1 & primary \\
Mortality, 365d & 365 & -0.0216 [-0.0523, +0.0098] & 0.177 & not corrected in this family & generating \\
\bottomrule
\end{tabular}
\par\footnotesize This family is corrected separately from the audit family (Additional file 1, Table S4) and the two are never pooled. No test carries an adjusted p in both.
\end{table}

% PANEL T6b
\begin{table}[t]
\centering
\small
\caption{The horizon ladder: trend statistics. Spearman rank correlation between horizon in days and the per-horizon quantity, with the interval taken on the correlation itself: patients are resampled once, every horizon is recomputed on that draw, and the correlation is recomputed from them.}
\label{tab:trends}
\begin{tabular}{llllll}
\toprule
Statistic & Role & Correlation [95\% CI] & Raw p & Holm p & Supported \\
\midrule
\texttt{acq\_advantage\_trend} & primary & -0.8000 [-1.0000, +0.2000] & 0.1416 & 0.708 & NOT SUPPORTED \\
\texttt{recovery\_ratio\_trend} & secondary & +0.2000 [-0.8000, +1.0000] & 0.5502 & 1 & NOT SUPPORTED \\
\texttt{acq\_advantage\_trend\_all\_six} & tertiary\_contains\_generating\_labels & -0.9429 [-1.0000, -0.1429] & 0.0026 & not corrected in this family & NOT SUPPORTED \\
\bottomrule
\end{tabular}
\par\footnotesize acq\_advantage\_trend\_all\_six contains the two observations that generated the hypothesis. It reaches a much stronger correlation with a small raw p, and it is reported for completeness only and never counted as support, whatever its value. That contrast is the entire argument for the restriction: the hypothesis looks convincing exactly when it is tested on the data that produced it.
\par\footnotesize An interval on a correlation over four points is coarse by construction, because such a correlation can take only a handful of values. That is a limitation of the pre-declared statistic and no amount of resampling fixes it.
\end{table}

% PANEL T7
\begin{table}[t]
\centering
\small
\caption{The noise floor, per arm and label, record level. Seed-to-seed spread of each cell's own five runs. The floor measured, 0.0087, is a crop-level spread of a 15-target macro and this project's contrasts are neither, so every contrast is judged against the spread of its own arms at its own label. Every per-label record-level spread of R6\_waveform exceeds that figure; several tabular arms fall well below it, which is why one scalar cannot serve them all.}
\label{tab:noise}
\begin{tabular}{llll}
\toprule
Arm & ICU admission, 24h & Mortality, 28d & Mortality, 365d \\
\midrule
\texttt{R1\_prevalence} & 0.0000 & 0.0000 & 0.0000 \\
\texttt{R2\_demo} & 0.0049 & 0.0024 & 0.0029 \\
\texttt{R3\_acqctx\_pre} & 0.0043 & 0.0207 & 0.0043 \\
\texttt{R3b\_acqctx\_window} & 0.0036 & 0.0047 & 0.0039 \\
\texttt{R4\_demo\_acq} & 0.0024 & 0.0045 & 0.0072 \\
\texttt{R5\_tabular} & 0.0016 & 0.0064 & 0.0029 \\
\texttt{R6\_waveform} & 0.0091 & 0.0190 & 0.0130 \\
\texttt{R7\_waveform\_demo\_acq} & 0.0020 & 0.0066 & 0.0042 \\
\bottomrule
\end{tabular}
\par\footnotesize A between-arm difference not clearly larger than the relevant cell here is reported as no difference.
\end{table}

% PANEL T8
\begin{table}[t]
\centering
\small
\caption{The waveform reproduction, and the only crop-level numbers here. The published macro is computed over crop-level rows, four per record , so the reproduction gap is computed against that protocol. Two published versions of the audited benchmark exist and they report different figures for this arm, so the gap is reported against both. The version of record is the primary comparison; the preprint is the figure this reproduction was judged against while it was the only one available. The record-level macro is the mean of the five per-seed macros and is reported beside it; the benchmark publishes nothing at that unit, so it has no gap. Neither number may enter Table 3's columns and neither is comparable to any tabular arm. The record-level macro here is a mean of per-seed macros, while Table 3's row 6 is one AUROC on the seed-pooled score, so the two differ slightly by construction.}
\label{tab:reproduction}
\begin{tabular}{lllllll}
\toprule
Arm & Unit & Published version & Local macro & Published [95\% CI] & Gap & Inside \\
\midrule
\texttt{R6\_waveform} & crop & version of record & 0.8474 & 0.8565 [0.8410, 0.8707] & -0.0091 & yes \\
\texttt{R6\_waveform} & crop & preprint & 0.8474 & 0.8279 [0.8047, 0.8511] & +0.0195 & yes \\
\texttt{R7\_waveform\_demo\_acq} & crop & version of record & 0.8797 & 0.8565 [0.8410, 0.8707] & +0.0232 & no \\
\texttt{R7\_waveform\_demo\_acq} & crop & preprint & 0.8797 & 0.8279 [0.8047, 0.8511] & +0.0518 & no \\
\texttt{R6\_waveform} & record & - & 0.8545 & none published at this unit & - & - \\
\texttt{R7\_waveform\_demo\_acq} & record & - & 0.8834 & none published at this unit & - & - \\
\bottomrule
\end{tabular}
\par\footnotesize The published interval is wide because six of the fifteen deterioration targets carry fewer than fifty positives in the test fold and the unweighted macro gives each the same weight as targets carrying over eight hundred. A pass is evidence the reproduction is not badly wrong, not that it is precisely right.
\par\footnotesize R7\_waveform\_demo\_acq is shown against the same published figures only for scale. It reproduces nothing, so falling outside is expected.
\par\footnotesize The local macro sits inside the published interval under both versions and on opposite sides of the published point estimate: below the version of record and above the preprint. No claim in this paper may rest on the direction of that gap.
\end{table}

% PANEL T9
\begin{table}[t]
\centering
\small
\caption{Calibration: expected calibration error under both declared schemes. ECE reads the probability scale; every other table here reads only the ranking, and the two can disagree completely. The reference is not zero: at this bin count and sample size a perfectly calibrated score still shows a non-zero ECE, so each cell carries the band such a score would produce, obtained by drawing labels from the arm's own predicted probabilities. An arm counts as calibrated only if it sits inside that band under both schemes, because the choice of binning is a choice.}
\label{tab:calibration}
\begin{tabular}{llllllll}
\toprule
Arm & Label & ECE, equal width [95\% CI] & ECE, equal mass [95\% CI] & Null band & Calibrated & Mean predicted & Observed \\
\midrule
\texttt{R1\_prevalence} & ICU admission, 24h & 0.0042 [0.0002, 0.0146] & 0.0042 [0.0002, 0.0146] & up to 0.0098 & yes & 0.1217 & 0.1260 \\
\texttt{R1\_prevalence} & Mortality, 28d & 0.0022 [0.0001, 0.0083] & 0.0022 [0.0001, 0.0083] & up to 0.0055 & yes & 0.0400 & 0.0422 \\
\texttt{R1\_prevalence} & Mortality, 365d & 0.0045 [0.0003, 0.0186] & 0.0045 [0.0003, 0.0186] & up to 0.0093 & yes & 0.1380 & 0.1425 \\
\texttt{R2\_demo} & ICU admission, 24h & 0.0082 [0.0024, 0.0187] & 0.0156 [0.0147, 0.0295] & up to 0.0179 & yes & 0.1218 & 0.1260 \\
\texttt{R2\_demo} & Mortality, 28d & 0.0023 [0.0006, 0.0084] & 0.0168 [0.0151, 0.0252] & up to 0.0109 & no & 0.0399 & 0.0422 \\
\texttt{R2\_demo} & Mortality, 365d & 0.0466 [0.0336, 0.0595] & 0.0495 [0.0420, 0.0654] & up to 0.0181 & no & 0.1374 & 0.1425 \\
\texttt{R3\_acqctx\_pre} & ICU admission, 24h & 0.0070 [0.0067, 0.0181] & 0.0085 [0.0090, 0.0205] & up to 0.0152 & yes & 0.1226 & 0.1260 \\
\texttt{R3\_acqctx\_pre} & Mortality, 28d & 0.0112 [0.0064, 0.0165] & 0.0196 [0.0161, 0.0255] & up to 0.0104 & no & 0.0403 & 0.0422 \\
\texttt{R3\_acqctx\_pre} & Mortality, 365d & 0.0194 [0.0114, 0.0309] & 0.0202 [0.0167, 0.0351] & up to 0.0182 & no & 0.1368 & 0.1425 \\
\texttt{R3b\_acqctx\_window} & ICU admission, 24h & 0.0170 [0.0135, 0.0269] & 0.0175 [0.0142, 0.0276] & up to 0.0161 & no & 0.1217 & 0.1260 \\
\texttt{R3b\_acqctx\_window} & Mortality, 28d & 0.0101 [0.0054, 0.0157] & 0.0102 [0.0082, 0.0171] & up to 0.0101 & no & 0.0404 & 0.0422 \\
\texttt{R3b\_acqctx\_window} & Mortality, 365d & 0.0105 [0.0071, 0.0237] & 0.0170 [0.0141, 0.0298] & up to 0.0175 & yes & 0.1366 & 0.1425 \\
\texttt{R4\_demo\_acq} & ICU admission, 24h & 0.0087 [0.0078, 0.0198] & 0.0094 [0.0086, 0.0204] & up to 0.0145 & yes & 0.1219 & 0.1260 \\
\texttt{R4\_demo\_acq} & Mortality, 28d & 0.0050 [0.0030, 0.0108] & 0.0093 [0.0066, 0.0152] & up to 0.0098 & yes & 0.0400 & 0.0422 \\
\texttt{R4\_demo\_acq} & Mortality, 365d & 0.0629 [0.0519, 0.0744] & 0.0643 [0.0534, 0.0754] & up to 0.0188 & no & 0.1364 & 0.1425 \\
\texttt{R5\_tabular} & ICU admission, 24h & 0.0103 [0.0099, 0.0204] & 0.0117 [0.0093, 0.0204] & up to 0.0127 & yes & 0.1206 & 0.1260 \\
\texttt{R5\_tabular} & Mortality, 28d & 0.0073 [0.0057, 0.0132] & 0.0077 [0.0066, 0.0143] & up to 0.0085 & yes & 0.0395 & 0.0422 \\
\texttt{R5\_tabular} & Mortality, 365d & 0.0190 [0.0138, 0.0291] & 0.0187 [0.0135, 0.0290] & up to 0.0147 & no & 0.1344 & 0.1425 \\
\texttt{R6\_waveform} & ICU admission, 24h & 0.1323 [0.1243, 0.1413] & 0.1347 [0.1259, 0.1428] & up to 0.0224 & no & 0.2502 & 0.1260 \\
\texttt{R6\_waveform} & Mortality, 28d & 0.1121 [0.1067, 0.1182] & 0.1120 [0.1059, 0.1178] & up to 0.0184 & no & 0.1542 & 0.0422 \\
\texttt{R6\_waveform} & Mortality, 365d & 0.1200 [0.1103, 0.1309] & 0.1202 [0.1099, 0.1306] & up to 0.0222 & no & 0.2551 & 0.1425 \\
\texttt{R7\_waveform\_demo\_acq} & ICU admission, 24h & 0.1303 [0.1234, 0.1385] & 0.1301 [0.1219, 0.1381] & up to 0.0209 & no & 0.2426 & 0.1260 \\
\texttt{R7\_waveform\_demo\_acq} & Mortality, 28d & 0.1102 [0.1047, 0.1158] & 0.1093 [0.1034, 0.1148] & up to 0.0179 & no & 0.1514 & 0.0422 \\
\texttt{R7\_waveform\_demo\_acq} & Mortality, 365d & 0.1291 [0.1194, 0.1397] & 0.1296 [0.1194, 0.1400] & up to 0.0226 & no & 0.2611 & 0.1425 \\
\bottomrule
\end{tabular}
\par\footnotesize Row 8 is absent by construction: it is row 6's scores under a restricted pair statistic, so it has no probability of its own and a calibration figure for it would be row 6's under another name.
\par\footnotesize The expected calibration error for R3\_acqctx\_pre at ICU admission, 24h under equal mass binning sits below the lower bound of its own percentile interval. Expected calibration error is positively biased, because binned absolute gaps cannot cancel and resampling noise can only add to them, so at a cell whose calibration error is near zero the interval can sit entirely above the value it was built around. The declared protocol is a percentile interval and it is reported as computed rather than swapped for a bias-corrected one; the null band, not the interval, is what answers whether the arm is calibrated.
\par\footnotesize Both waveform arms over-predict at every label, by a factor running from about 1.8 to about 3.7 of the observed rate. The largest factor is at 28 days, where the observed rate is lowest; both columns it is computed from are in this table. They were fitted under focal BCE and their score is a mean over four crops and then over five seeds, and each of those moves the probability scale.
\end{table}

% PANEL T10
\begin{table}[t]
\centering
\small
\caption{Net benefit at the pre-declared threshold. Threshold 0.10, fixed from the clinical framing and recorded before any curve was computed, and quoted at 24-hour ICU admission alone because it is the only reported label with an action available at emergency-department disposition. Treating everyone scores 0.0288 here and treating no one scores exactly zero; the prevalence is 0.1260, which no rule can exceed. The comparison against treat-all is a paired clustered bootstrap on the difference, since treat-all also varies with the resampled prevalence.}
\label{tab:net_benefit}
\begin{tabular}{llllll}
\toprule
Arm & Net benefit [95\% CI] & vs treat-all [95\% CI] & Beats treat-all & Fraction flagged & Mean predicted \\
\midrule
\texttt{R5\_tabular} & 0.0897 [0.0802, 0.0993] & +0.0609 [+0.0560, +0.0654] & yes & 0.2397 & 0.1206 \\
\texttt{R4\_demo\_acq} & 0.0671 [0.0580, 0.0766] & +0.0383 [+0.0334, +0.0430] & yes & 0.3902 & 0.1219 \\
\texttt{R3b\_acqctx\_window} & 0.0662 [0.0570, 0.0759] & +0.0374 [+0.0327, +0.0420] & yes & 0.4122 & 0.1217 \\
\texttt{R3\_acqctx\_pre} & 0.0653 [0.0564, 0.0748] & +0.0365 [+0.0317, +0.0411] & yes & 0.4077 & 0.1226 \\
\texttt{R7\_waveform\_demo\_acq} & 0.0367 [0.0254, 0.0482] & +0.0079 [+0.0069, +0.0088] & yes & 0.9277 & 0.2426 \\
\texttt{R2\_demo} & 0.0351 [0.0243, 0.0461] & +0.0062 [+0.0022, +0.0101] & yes & 0.7987 & 0.1218 \\
\texttt{R6\_waveform} & 0.0334 [0.0220, 0.0449] & +0.0045 [+0.0038, +0.0053] & yes & 0.9575 & 0.2502 \\
\texttt{R1\_prevalence} & 0.0288 [0.0174, 0.0403] & +0.0000 [-0.0000, +0.0000] & no & 1.0000 & 0.1217 \\
\bottomrule
\end{tabular}
\par\footnotesize Net benefit reads the probability scale rather than the ranking, so this ordering is not a discrimination ordering: an arm that ranks cases well loses at a fixed threshold when its probabilities are inflated. No recalibration was performed; only test-fold predictions are stored, and refitting on them would be circular.
\end{table}

% Footnote, generated with the tables:
% - Resampling: patient-clustered percentile bootstrap, 10,000 resamples,
%   resampling unit subject\_id. Contrasts are paired on the same
%   draw; two marginal intervals are never compared.
% - Unit: record level, one row per ECG, for Tables 1, 3, 6a, 6b and 7.
%   Tables 4a and 4b are a within-stratum concordance over restricted pairs
%   and are a different estimand, which is why row 8 is not a row of Table 2.
%   Table 9 is the only place crop-level numbers appear.
% - Multiplicity: two families, corrected separately and never pooled.
%   acq\_context\_audit at its declared size m = 15;
%   horizon\_ladder at its declared size m = 6. Both use Holm.
% - Seeds: five per arm. Scores are averaged across seeds and one AUROC is
%   taken on the pooled score; the per-seed spread is Table 8.
% - Noise floor: per (arm, label) at record level, not a single scalar.
% - Row-level predictions live outside the repository.
